\section{Notación básica}
\subsection{Conjuntos}
Sea $\Gamma=\set{k_1,\dots,k_l}\subseteq\N_n$ un conjunto con cardinalidad $|\Gamma|=l$ y $\gamma=(\gamma_{k})_{k=1}^n\in\B^n$; denotamos por $\ext{\gamma}{\Gamma}=(\ext{\gamma}{\Gamma}_k)_{k=1}^n$ el elemento de $\mathbb{Z}_2^n$ definido por:
\begin{equation*}
  \ext{\gamma}{\Gamma}_k=
  \begin{cases}
    \gamma_k & \text{ si }k\in\Gamma    \\
    0        & \text{ si }k\notin\Gamma
  \end{cases}\,.
\end{equation*}

Por otro lado, si $\gamma\in\mathbb{Z}^n$, llamaremos $\con{\gamma}{\Gamma}$ al elemento de $\mathbb{Z}^l$ definido por $\con{\gamma}{\Gamma}=(\gamma_{k_m})_{m=1}^l$.

Observemos que para todo elemento $\gamma\in\mathbb{Z}^l$ se cumple que $\con{\ext{\gamma}{\Gamma}}{\Gamma}=\gamma$ mientras que en general si $\gamma\in\mathbb{Z}^n$ entonces $\ext{\con{\gamma}{\Gamma}}{\Gamma}\neq\gamma$.

También vamos a considerar que los elementos $(0,\dots,0)\in\B^l$ y $(1,\dots,1)\in\B^l$ lo denotaremos simplemente por $0$ y $1$.

\subsection{Qubits}
Como se introduce en~\cite{diracBraket}, los qubits se escriben normalmente en la llamada notación bra-ket.
Cada qubit se representa como un ket, de la forma $\ket{a}$, que denota un vector en $\Q$.

Por ejemplo, los estados de la base computacional se escriben como $\ket{0}$ y $\ket{1}$, llamados ket 0 y ket 1, respectivamente (véase~\cite{NielsenChuang}).
Así, cualquier qubit puede describirse como una combinación lineal de estos dos estados de la base: $\ket{a} =\alpha_0\ket{0} +\alpha_1 \ket{1}$, donde $\alpha_0,  \alpha_1 \in \mathbb{C}$ con $ |\alpha_0|^{2}+|\alpha_1|^{2}=1$.

\begin{notation}
  Para cualquier qubit $\ket{a}$, se usa la siguiente notación:
  \[a_{0}=\braket{0}{a}\text{ y }a_{1}=\braket{1}{a},\]
\end{notation}
por lo que el qubit $\ket{a}$ puede escribirse como:
\begin{equation}
  \label{eq:qubit-coordenadas}
  \ket{a}=a_0\ket{0}+a_1\ket{1}
\end{equation}

Para otras bases, como por ejemplo $\set{\ket{+}, \ket{-}}$, donde $\ket{+}=\frac{\ket{0}+\ket{1}}{\sqrt{2}}$ y $\ket{-}=\frac{\ket{0}-\ket{1}}{\sqrt{2}}$, un qubit $\ket{a}$ puede describirse como:
\begin{equation*}
  \ket{a}=a_+\ket{+}+a_-\ket{-}
\end{equation*}
